\section{Checks on the curve and the noise}\label{app:curve-checks}

Four checks bear on the curve of Sec.~\ref{sec:curve} and the noise of Sec.~\ref{sec:measure}.
They use development runs on the demonstration set only, and the curve was built after most of what they test. The tables call the ceiling of Sec.~\ref{sec:curve} the plateau, and a development run a replica.

\paragraph{Runs needed to see a serving gain.}
At the curve's elasticity of 0.31, +11\% tokens is worth 4.2 of 124 levels. Seeing it takes about 32 runs per arm at the SD of eight identical runs made on one day (5.95) and 61 at the pairs' SD (8.25), both two-sided, as every detection number in the paper is (Table~\ref{tab:detect}). % C4 V3-X13 V3-W12
Our promotion tests use two runs per arm.

\begin{table}[!htb]
\centering\footnotesize
\caption{Runs per arm needed to see a token gain on the replica (80\% power, two-sided 5\% level, normal approximation, shown before rounding; round up for the stated power). The gain is the curve's prediction (elasticity 0.31 on levels) out of 124 levels. SD 5.95: eight byte-identical runs made on one day (six started together); SD 8.25: the pooled SD of the six identical pairs. Every detection number in the paper is two-sided at an identical-run SD.}\label{tab:detect}
\setlength{\tabcolsep}{4pt}
\begin{tabular}{@{}rrrr@{}}
\toprule
Extra tokens & Gain (levels) & SD 5.95 & SD 8.25 \\
\midrule
+5\% & 1.9 & 153.1 & 294.2 \\ % C4 V3-W12
+11\% & 4.2 & 31.6 & 60.8 \\ % C4 V3-W12
+26\% & 9.9 & 5.7 & 10.9 \\ % C4 V3-W12
\bottomrule
\end{tabular}
\end{table}

\paragraph{Past serving steps.}
Table~\ref{tab:steps} sets twelve earlier steps that moved tokens per second against the curve. Fitted together they give an elasticity of 0.45 (95\% interval 0.13 to 0.77 at the pooled SD of 8.25; 0.29 to 0.61 at 4.1), next to the curve's 0.31. % N1.3
The pruning step carries the fit: $+19.5$ levels at $+11.2\%$ tokens against 3.3 predicted. Without it the estimate is 0.02 ($-0.48$ to 0.53), and the ten steps after pruning lean below the curve at the shipped family's SD. % N1.4
One step's noise (8 to 12 levels) swamps the curve's prediction (4.5 levels at most), so the steps neither confirm nor refute the curve. We built it after all of them, and none of the 29 promotion-test records of this stack holds a level prediction from tokens. % N1.5 N1.8

% Generated by paper/v2/analysis/new_token_steps.py (N1). Do not edit by hand.
\begin{table}[!htb]
\caption{Steps that moved tokens per second, each an arm against its control on the development replica (25 public games, not held out). Tokens: generated tokens, arm over control. Levels: control to arm, sums over 25 games (means where several runs exist). Curve: $0.31\times$ relative token change $\times$ control levels (Sec.~\ref{sec:curve}). Runs: arm/control replicas; the control is every archived run of the preceding configuration. Output kept: yes, outputs given the prompt are unchanged; no, the model or its prompt changes; fix, a correctness repair; part, the refit is lossless but the pair also changes the prompt format. Steps share runs, and the curve was built after all of them. Fitted through the origin, the twelve serving steps give an elasticity of 0.45 (95\% interval 0.13 to 0.77 at the pooled run SD of 8.25 levels, 0.29 to 0.61 at 4.1); without the pruning step, 0.02 ($-0.48$ to 0.53; $-0.23$ to 0.28). Twelve more pairs that change the agent or the scheduler are in the evidence file.}\label{tab:steps} % N1 C3 C4 C9 B1
\centering\footnotesize
\setlength{\tabcolsep}{3pt}
\begin{tabular}{@{}>{\raggedright\arraybackslash}p{0.30\linewidth}lrrrrr@{}}
\toprule
Step & Output kept & Runs & Tokens (\%) & Levels & $\Delta$ levels & Curve \\
\midrule
12 streams, no memory freed & yes & 1/2 & $-15.6$ & 94.5 to 90 & $-4.5$ & $-4.5$ \\ % N1 S01
REAP keep-448 experts, 12 streams & no & 2/2 & $+11.2$ & 94.5 to 114 & $+19.5$ & $+3.3$ \\ % N1 S02
14 streams, 70 Mamba slots & yes & 2/2 & $+11.6$ & 114 to 114 & $0.0$ & $+4.1$ \\ % N1 S03
Host-RAM KV tier at 14 streams & yes & 1/2 & $-0.8$ & 114 to 107 & $-7.0$ & $-0.3$ \\ % N1 S04
16 streams, host KV tier & yes & 1/2 & $+6.8$ & 114 to 115 & $+1.0$ & $+2.4$ \\ % N1 S05
Relaxed MTP acceptance (0.5) & no & 1/2 & $+6.4$ & 114 to 113 & $-1.0$ & $+2.2$ \\ % N1 S06
16 streams, host KV tier (D-prime stack) & yes & 1/2 & $+4.6$ & 112 to 104 & $-8.0$ & $+1.6$ \\ % N1 S08
QSA index fix, async M-RoPE upload & fix & 1/2 & $+3.3$ & 126 to 119 & $-7.0$ & $+1.3$ \\ % N1 S12
Image hashes off scheduler, toolfast, watchdog & yes & 1/1 & $+5.7$ & 119 to 128 & $+9.0$ & $+2.1$ \\ % N1 S15
Drafter refit (D3), prompt format & part & 2/1 & $+5.0$ & 128 to 125.5 & $-2.5$ & $+2.0$ \\ % N1 S16
Fourth draft step & yes & 1/2 & $-2.5$ & 125.5 to 120 & $-5.5$ & $-1.0$ \\ % N1 S17
32K draft vocabulary map & yes & 1/2 & $+1.0$ & 125.5 to 126 & $+0.5$ & $+0.4$ \\ % N1 S18
\midrule
Yield at 8,192 tokens (the yield pair) & no & 2/2 & $-0.2$ & 111 to 126 & $+15.0$ & $-0.1$ \\ % N1 S11
\bottomrule
\end{tabular}
\end{table}


\paragraph{Plateau against pool size.}
The plateau estimate rises as more runs are pooled, because a level that no pooled attempt cleared stays stuck and pooling unsticks some (Table~\ref{tab:plateau}). % N2
The share of it that a replica run earns is 94\% at eight runs, 92\% at ten and 82\% to 85\% when 21 to 34 runs are pooled, and extrapolates to 85\% to 91\%: about 90\%. The 94\% is the eight-run value, not an extrapolation. % N2 C9
A bootstrap over the eight runs gives a simulated share of 94.7\% (93.4\% to 96.0\%), so pool size moves it more than run sampling does. The elasticity stays at 0.31 to 0.33 from two to ten pooled runs, so the 4-level price does not move. % N2 C3
The plateau used here is the expected score at unlimited budget with each level redrawn from its recorded attempts; counting every level that any attempt cleared would give 70.9 on the eight-run pool. On the older September stack the run share is 59\% at 31 runs, so the near-plateau result is specific to our stack. % N3-11 N2

% Generated by paper/v2/analysis/new_plateau_pool.py (N2). Do not edit by hand.
\begin{table}[!htb]
\caption{Plateau $P$ against the number $k$ of pooled replica runs, and the share of it that a replica run earns. $P$ is the expected mean game score if every level the pooled attempts can clear is cleared (Sec.~\ref{sec:curve}), averaged over all subsets of $k$ runs (or 200 random ones). $S(1)/P$: simulated score at the replica budget over $P$, in percent, on the same 110 games per seed. Run share: the pooled runs' own mean score over $P$, in percent, without the simulator. $k=10$ adds two later runs that were not in the paper's pool. Limits: fits of a hyperbolic and an exponential form to the mean curves and transfers from the control and all-replica curves; the stress case does not scale the transfer by the pool's slower growth; the cap treats every level that some eight-run replica cleared as always clearable. Development replica (25 public games), not held out.}\label{tab:plateau} % N2 C9
\centering\footnotesize
\setlength{\tabcolsep}{4pt}
\begin{tabular}{@{}lrrrr@{}}
\toprule
Pool & $k$ & $P$ & $S(1)/P$ (\%) & Run share (\%) \\
\midrule
Eight-run pool, plus two later runs & 1 & 56.8 & 96 & 100 \\ % N2 C9 pin k=1
 & 2 & 57.9 & 95 & 98 \\ % N2 C9 pin k=2
 & 4 & 59.1 & 95 & 96 \\ % N2 C9 pin k=4
 & 8 & 60.2 & 94 & 94 \\ % N2 C9 pin k=8
 & 10 & 61.4 & 94 & 92 \\ % N2 C9 pin k=10
\addlinespace[2pt]
Yield-2,048 controls (21 runs) & 1 & 49.6 & 97 & 100 \\ % N2 C9 ctrl21 k=1
 & 2 & 51.0 & 95 & 98 \\ % N2 C9 ctrl21 k=2
 & 4 & 52.7 & 94 & 94 \\ % N2 C9 ctrl21 k=4
 & 8 & 54.8 & 92 & 90 \\ % N2 C9 ctrl21 k=8
 & 21 & 58.8 & 88 & 85 \\ % N2 C9 ctrl21 k=21
\addlinespace[2pt]
All 34 replicas & 1 & 50.9 & n/a & 100 \\ % N2 C9 fz34 k=1
 & 4 & 54.6 & n/a & 93 \\ % N2 C9 fz34 k=4
 & 8 & 56.9 & n/a & 90 \\ % N2 C9 fz34 k=8
 & 21 & 60.4 & n/a & 84 \\ % N2 C9 fz34 k=21
 & 34 & 61.8 & n/a & 82 \\ % N2 C9 fz34 k=34
\addlinespace[2pt]
September stack (27B, 31 runs) & 1 & 5.2 & n/a & 100 \\ % N2 C9 sept27 k=1
 & 4 & 5.9 & n/a & 87 \\ % N2 C9 sept27 k=4
 & 8 & 6.6 & n/a & 78 \\ % N2 C9 sept27 k=8
 & 31 & 8.8 & n/a & 59 \\ % N2 C9 sept27 k=31
\addlinespace[2pt]
Eight-run limit: five estimates & $\infty$ & 62.1 to 66.4 & 92 to 94 & 85 to 91 \\ % N2 limits
 & & median 64.6 & n/a & 88 \\ % N2 limits
Eight-run limit: stress case & $\infty$ & 68.6 to 69.4 & n/a & 82 to 83 \\ % N2 limits
Eight-run cap: no new level & $\infty$ & 70.6 & n/a & 80 \\ % N2 cap
\bottomrule
\end{tabular}
\end{table}




\paragraph{A noise model and the stuck mass.}
A model that clears every level independently at its recorded rate, under the replica's token budget, predicts run SDs of 5.2, 6.4 and 5.7 levels, against 3.9, 5.5 and 8.25 observed (Table~\ref{tab:noisemodel}); the standard error of each prediction is 0.3 to 0.8. % N3-3
It matches the stable groups and underpredicts the identical pairs by a factor 1.44 in SD (the chance of an SD this high under the model is 0.05). The gap is the pairs' common first-minus-second shift of 9.2 levels, which holds 62\% of their sum of squares; around it the pairs' SD is 5.6 against the model's 5.7. Tokens per second and slot contention do not explain it. % N3-7
Completion explains 84\% to 90\% of within-game score variance on our replica, as the forum post reports (84.7\%, by our reading of its definition \citep{borro19802026forum732932}), but the two largest games hold 27\% to 32\% of it, against 65\% there. % N3-9
Of the points the nine yield runs lose to unreached levels, 57\% sit on levels that none of 30 runs cleared; the table splits the rest by tokens spent. % N3-12

% Generated by paper/v2/analysis/new_noise_model.py (--tex-only rebuilds it from out/new_noise_model.json). Label owner: N3.
\begin{table}[!htb]
\centering\footnotesize\setlength{\tabcolsep}{4pt}
\caption{A noise model and the stuck mass, on the development replica (25 public games, not held out).
\emph{Top:} SD across runs of total levels and of mean game score, observed and predicted from recorded level attempts.
\emph{Binary} clears every level independently at its recorded rate; \emph{Budget} also plays the recorded costs on 14 shared streams in 7,200 s (SE: bootstrap, then jackknife, over runs).
The yield pool is eight runs: seven of the nine yield runs compared in Sec.~\ref{sec:yield} and the run without a token yield.
\emph{Middle:} a level is never cleared when no attempt in the pool cleared it; points are the share of a game's 100 held by such levels, averaged over 25 games, as a share of the points the pool's own runs lose to unreached levels.
\emph{Bottom:} points per game lost to unreached levels, by where the run stalled, against all 30 runs.}
\label{tab:noisemodel}
\begin{tabular}{@{}p{3.1cm}rrrrrr@{}}
\toprule
 & & \multicolumn{3}{c}{SD of total levels} & \multicolumn{2}{c}{SD of mean score} \\
\cmidrule(lr){3-5}\cmidrule(l){6-7}
Compared runs & Attempts & Observed (95\% CI) & Binary & Budget (SE) & Obs. & Budget \\
\midrule
Nine yield-8,192 runs & 1,114 & 3.9 (2.6 to 7.5) & 4.6 & 5.2 (0.4, 0.6) & 2.7 & 3.1 \\ % N3-3
21 yield-2,048 controls & 1,862 & 5.5 (4.2 to 8.0) & 6.4 & 6.4 (0.3, 0.8) & 3.6 & 4.0 \\ % N3-3
Six identical pairs, within pair & 1,588 & 8.25 (5.3 to 18.2) & 5.8 & 5.7 (0.3, 0.3) & 5.0 & 3.8 \\ % N3-3
\bottomrule
\end{tabular}

\medskip
\begin{tabular}{@{}p{3.1cm}rrrrr@{}}
\toprule
Pool & Attempts & Levels never cleared (of 183) & Games & Points held & Share \\
\midrule
Eight runs (seven at 8,192, one with no token yield) & 1,114 & 35 & 13 & 26.1 & 64\% \\ % N3-11
Twelve plain runs & 1,862 & 40 & 13 & 29.0 & 57\% \\ % N3-11
All 30 runs & 3,969 & 31 & 12 & 23.2 & 50\% \\ % N3-11
\bottomrule
\end{tabular}

\medskip
\begin{tabular}{@{}p{5.2cm}rrrr@{}}
\toprule
 & \multicolumn{2}{c}{Nine yield runs} & \multicolumn{2}{c}{21 controls} \\
\cmidrule(lr){2-3}\cmidrule(l){4-5}
Unreached points by where the run stalled & Points & Share & Points & Share \\
\midrule
No run cleared the level & 23.2 & 57\% & 23.2 & 48\% \\ % N3-12
Stalled with at least the tokens of the longest clear of that level & 9.9 & 24\% & 9.8 & 20\% \\ % N3-12
Stalled with fewer tokens than that & 7.4 & 18\% & 15.4 & 32\% \\ % N3-12
All unreached levels & 40.5 & 100\% & 48.3 & 100\% \\ % N3-12
\bottomrule
\end{tabular}
\end{table}

